\centerline{\bf \Huge Разнобой}


\problem{1}
Дана функция $f$, определённая на множестве действительных чисел и принимающая действительные значения. Известно, что для любых $x$ и $y$ таких, что $x>y$, верно неравенство $(f(x))^2 \leqslant f(y)$. Докажите, что множество значений функции содержится в промежутке $[0,1]$.

\problem{2}
Существует ли бесконечная последовательность натуральных чисел такая, что для любого натурального $k$ сумма любых $k$ идущих подряд членов этой последовательности делится на $k+1$?

\problem{3}
Квадрат разбит на $n^2 \geqslant 4$ прямоугольников $2(n-1)$ прямыми, из которых $n-1$ параллельны одной стороне квадрата, а остальные $n-1$ - другой. Докажите, что можно выбрать $2 n$ прямоугольников разбиения таким образом, что для любых двух выбранных прямоугольников один из них можно поместить в другой (возможно, предварительно повернув).

\problem{4}
Пусть $P(x)$ - многочлен степени $n \geqslant 2$ с неотрицательными коэффициентами, а $a, b$ и $c$ - длины сторон некоторого остроугольного треугольника. Докажите, что числа $\sqrt[n]{P(a)}, \sqrt[n]{P(b)}$ и $\sqrt[n]{P(c)}$ также являются длинами сторон некоторого остроугольного треугольника.

\problem{5}
Даны натуральные числа $a$ и $b$. Докажите, что существует бесконечно много натуральных $n$ таких, что число $a^n+1$ не делится на $n^b+1$.

\problem{6}
Паша и Вова играют в следующую игру, делая ходы по очереди. Начинает Паша. Изначально перед мальчиками лежит большой кусок пластилина. За один ход Паша может разрезать любой из имеющихся кусков пластилина на три части (не обязательно равные). Вова своим ходом выбирает два куска и слепляет их вместе. Паша побеждает, если в некоторый момент среди имеющихся кусков пластилина окажется 100 кусков одинаковой массы. Может ли Вова помешать Паше победить?